%==============================================================================
% Abstract - Tóm tắt luận văn
%==============================================================================

\documentclass[12pt]{article}
\usepackage[utf8]{vietnam}
\usepackage{times}

\begin{document}

\title{\textbf{Nghiên cứu và đề xuất phương pháp sinh câu hỏi trắc nghiệm tự động dựa trên mô hình ngôn ngữ lớn kết hợp hệ thống đa tác tử}}
\author{Nguyễn Văn A}
\date{2024}
\maketitle

%------------------------------------------------------------------------------

\section*{Tóm tắt}

Trong bối cảnh giáo dục hiện đại, nhu cầu tự động hóa quá trình tạo ngân hàng câu hỏi trắc nghiệm ngày càng tăng cao nhằm giảm tải giáo viên và đảm bảo tính nhất quán trong đánh giá. Luận văn này nghiên cứu và đề xuất phương pháp sinh câu hỏi trắc nghiệm tự động dựa trên mô hình ngôn ngữ lớn (Large Language Models - LLMs) kết hợp với kiến trúc hệ thống đa tác tử (Multi-Agent Systems).

Để đảm bảo câu hỏi được phân loại theo mức độ nhận thức phù hợp với mục tiêu học tập, nghiên cứu áp dụng phân loại Bloom (Bloom's Taxonomy) với ba cấp độ: Truy xuất (Retrieval), Suy luận và tổng hợp (Inference \& Synthesis), và Lập luận phản biện (Critical Reasoning). Phương pháp đề xuất sử dụng kiến trúc Multi-Agent dựa trên framework LangGraph, trong đó các tác tử chuyên biệt phụ trách từng nhiệm vụ: tác tử phân tích văn bản đầu vào, các tác tử sinh câu hỏi theo từng cấp độ Bloom, và tác tử tổng hợp kết quả.

Ngoài ra, luận văn xây dựng quy trình OCR (Optical Character Recognition) hoàn chỉnh bao gồm: sinh ảnh sạch từ văn bản, tăng cường ảnh sử dụng thư viện Augraphy để mô phỏng điều kiện scanning thực tế, và nhận dạng văn bản bằng cả phương pháp truyền thống (EasyOCR, PaddleOCR) và mô hình thị giác-ngôn ngữ (GPT-4V, Claude Vision). Khung đánh giá đa chiều được đề xuất với bốn tiêu chí: đánh giá vi mô (Solvability, Distractor Quality, Alignment), đánh giá vĩ mô (LLM-as-a-Judge), đánh giá định dạng đầu ra, và đánh giá chất lượng OCR qua các chỉ số CER (Character Error Rate) và WER (Word Error Rate).

Thực nghiệm được tiến hành trên tập dữ liệu RACE (Reading Comprehension) với việc so sánh giữa các mô hình LLMs khác nhau (GPT-4, Claude-3, LLaMA-3) và giữa hai phương pháp sinh câu hỏi: đơn prompt (Single-Prompt) và đa tác tử (Multi-Agent). Kết quả cho thấy phương pháp Multi-Agent với thiết kế prompt riêng biệt cho từng cấp độ Bloom đạt được chất lượng câu hỏi tốt hơn, đặc biệt trong việc tạo câu hỏi yêu cầu suy luận và lập luận phản biện. Nghiên cứu cũng phân tích chi tiết các yếu tố ảnh hưởng đến chất lượng đầu ra và đề xuất các hướng phát triển trong tương lai.

\textbf{Từ khóa:} Sinh câu hỏi trắc nghiệm tự động, Mô hình ngôn ngữ lớn, Hệ thống đa tác tử, Phân loại Bloom, Nhận dạng ký tự quang học, LangGraph.

%------------------------------------------------------------------------------
\pagebreak

\section*{Abstract}

In the context of modern education, the demand for automated test question generation is increasing to reduce teacher workload and ensure consistency in assessment. This thesis proposes an automatic multiple-choice question generation method based on Large Language Models (LLMs) combined with Multi-Agent Systems architecture.

To ensure questions are categorized according to appropriate cognitive levels aligned with learning objectives, this research applies Bloom's Taxonomy with three levels: Retrieval, Inference \& Synthesis, and Critical Reasoning. The proposed method utilizes a Multi-Agent architecture based on the LangGraph framework, with specialized agents handling distinct tasks: an analyzer agent for processing input text, generator agents for each Bloom level, and a merge agent for synthesizing results.

Additionally, this thesis develops a complete Optical Character Recognition (OCR) pipeline including: clean image generation from text, image augmentation using Augraphy library to simulate real scanning conditions, and text recognition using both traditional methods (EasyOCR, PaddleOCR) and Vision-Language Models (GPT-4V, Claude Vision). A multi-dimensional evaluation framework is proposed with four criteria: micro evaluation (Solvability, Distractor Quality, Alignment), macro evaluation (LLM-as-a-Judge), output format evaluation, and OCR quality evaluation using Character Error Rate (CER) and Word Error Rate (WER) metrics.

Experiments are conducted on the RACE (Reading Comprehension) dataset, comparing different LLM models (GPT-4, Claude-3, LLaMA-3) and two question generation methods: Single-Prompt and Multi-Agent. Results show that the Multi-Agent method with level-specific prompt designs achieves better question quality, particularly for inference and critical reasoning questions. The study also analyzes factors affecting output quality and proposes future research directions.

\textbf{Keywords:} Automatic Multiple-Choice Question Generation, Large Language Models, Multi-Agent Systems, Bloom's Taxonomy, Optical Character Recognition, LangGraph.

\end{document}
